\documentclass[11pt]{article}
\usepackage[margin=1in]{geometry}
\usepackage{amsmath, amssymb}
\usepackage{fancyhdr}
\usepackage{enumitem}
\usepackage{xcolor}
\usepackage{graphicx}
\usepackage{xparse}       % For defining commands with multiple arguments
\usepackage{float}
\usepackage{subcaption}
\usepackage{booktabs}         % For \midrule, \toprule, \bottomrule in tables

% Custom problem command
\NewDocumentCommand{\problem}{m m}{%
  \section*{Problem #1}           % 2. Section heading with problem number
  \input{#2}
  \vspace{1em}\par\noindent\hrule % 1. Horizontal line after the problem
}
\newcommand{\question}[1]{\fbox{\parbox{\textwidth}{#1}}\vspace{1em}}



% Title info
\newcommand{\hmwkTitle}{Problem Set 5}
\newcommand{\hmwkDueDate}{\today}
\newcommand{\hmwkClass}{EC 708 - Econometrics 1}
\newcommand{\hmwkAuthorName}{Kartik Patekar U71448411}


% Useful math symbols
\newcommand{\R}{\mathbb{R}}
\newcommand{\N}{\mathbb{N}}


% Header/Footer
\pagestyle{fancy}
\fancyhf{}
\rhead{\hmwkAuthorName}
\lhead{\hmwkTitle}
\cfoot{\thepage}


% Document
\begin{document}

% no indentation
\setlength{\parindent}{0pt}

% Title
\begin{center}
  \textbf{\Large \hmwkClass}\\
  \textbf{\hmwkTitle}\\
  \hmwkAuthorName \\
  \hmwkDueDate
\end{center}

% -------------------
\problem{1}{Exercises/ex1.tex}


% -------------------


% \begin{figure}[h!]
%   \centering
%   \includegraphics[width=0.8\textwidth]{example-image.png}
%   \caption{A sample image}
%   \label{fig:sample-image}
% \end{figure}

% -------------------
\end{document}
